\chapter{Which Questions Still Have No Answers?}

\begin{kaichang}
If you visit the seaside, try something simple: bring home a bottle of seawater and place it on your desk. It looks clear and clean, little different from ordinary bottled mineral water.

Now consider three facts about that bottle. First, one milliliter of seawater---approximately twenty drops---typically contains more than a million bacteria and ten million virus particles, a population denser than a major city's. Second, scientists still cannot name most of these residents. We know they exist and can read their ``signatures'' in DNA sequences, but have never cultured them in a laboratory and examined them clearly. Third, and most important, the bottle itself is the scene of an unsolved case. Approximately four billion years ago, in water like this or somewhere near similar water, ``ordinary chemicals'' made the leap to ``a self-replicating system.'' Exactly how that happened remains unknown today.

Across eighty-six chapters, this book has explained how atoms build molecules, molecules form cells, and information is stored and rewritten in DNA. You might now think science has understood every link in this chain. The final chapter has the opposite task: to visit the lamps along that chain that remain unlit.
\end{kaichang}

\textbf{Translator safety note:} A clear sample is not necessarily safe to drink or handle. Collect only with adult supervision from a safe shore, avoid suspect blooms, keep samples labeled and separate from drinks, and wash hands after handling. Follow \href{https://cdc.gov/harmful-algal-blooms/prevention/index.html}{CDC water and algal-bloom precautions}.

\section{Different Kinds of ``We Don't Know''}

First, a scientist's ``I don't know'' is not the same as a student's ``I don't know'' in an examination.

In an exam, it often means ``The answer is in the book, but I forgot it.'' At science's frontier, there are at least three kinds. The first is \textbf{not yet investigated}: nobody has sampled a particular deep-ocean trench to discover its species; nobody has looked. The second is \textbf{under investigation, with candidate answers}: competing hypotheses each explain some evidence, and researchers are designing experiments to distinguish them. Most questions in this chapter belong here. The third is deeper: \textbf{we are not yet sure which concepts to use in asking}. For consciousness, for example, there is no agreed method for measuring whether a system has it; the question itself is still being refined.

You have encountered these kinds in everyday life. A doctor saying ``The cause of this symptom is unclear'' may mean it has not yet been found, the first kind, or that several possibilities are being ruled out, the second. A forecast of a 60\% chance of rain in three days is not laziness: the atmosphere is extremely sensitive to initial conditions, so tiny measurement differences can produce completely different outcomes days later. This ``probabilities are all we can give in principle'' kind of uncertainty will return later.

Three criteria help decide whether an unknown is a good scientific question: clear boundaries, so we know which parts are already understood; candidate explanations, distinguishable guesses rather than a blank; and a next step, an experiment or observation that can actually be made. With this yardstick, let us visit five major questions still open today.

\section{How Exactly Did Life Begin?}

Turn the camera back approximately four billion years. Earth's oceans held no fish or algae, and perhaps almost no oxygen, just water, dissolved salts, simple gases, and energy supplied by volcanoes and lightning. Every bacterium in today's bottle descends from the transformation that followed.

What do we already know? Both ends of the chain are fairly clear. At the beginning, simple inorganic and organic molecules can spontaneously assemble into life's ``parts''---amino acids, nucleotides, and fatty acids---under suitable conditions. Experiments and meteorites repeatedly support this, as the evidence story below explains. At the other end, all living organisms share the same genetic code and core molecular machinery, indicating common ancestry; Chapters 45 and 56 discussed this evidence. The ribosome, Chapter 69's protein-translation machine, has an RNA core, like a living fossil suggesting that an ``RNA world'' may once have existed, with RNA serving both as information carrier and catalyst.

The middle is missing. How many steps were there between ``reactive small molecules'' and ``a system that copies itself, makes mistakes, and undergoes natural selection''? Did it happen in mineral micropores at deep-sea vents, in shallow pools that alternately dried and refilled near volcanoes, or within ice? Did lipid vesicles (Chapter 47) or RNA come first? Each question has serious candidate answers and dedicated experiments, but none can claim victory. The honest answer is: \textbf{the broad picture of life's origin is credible; the specific route is unknown}.

Two leading settings each have strengths and weaknesses. Deep-sea-vent proponents point to mineral micropores as countless tiny reaction chambers, iron--sulfur minerals on their walls as catalysts, and the temperature difference between hot vent water and cold seawater as a continuous energy flow. Many microbes in today's depths indeed live on such chemical energy. Critics respond that seawater is too abundant and salty for components such as nucleotides to concentrate enough to collide and join into chains. Shallow pools offer the opposite advantage: repeated drying and wetting concentrates, dries, and redisperses molecules, and wet--dry cycles have genuinely produced long molecular chains in laboratories. But were there enough stable pools on land four billion years ago? Would strong ultraviolet radiation quickly destroy newly formed chains? Each camp can experiment and be challenged, and neither has yet been disproved. This is what ``under investigation, with candidate answers'' really looks like.

\begin{zhengju}
In 1953, a University of Chicago graduate student named Miller, supported by his adviser Urey, assembled a glass apparatus. One flask held water representing the ocean, heated to produce ``ocean vapor.'' Another large flask contained methane (\chem{CH_4}), ammonia (\chem{NH_3}), and hydrogen (\chem{H_2}), representing the early atmosphere. Electrical sparks between them simulated lightning. After a week of circulation, several amino acids appeared in the water: protein components had formed on their own in an ``inorganic soup.'' The sensational result first demonstrated that ``the first step from chemistry to life'' could be tested in a laboratory.

The second half of the story matters just as much. Later studies reassessed the early atmosphere. Volcanic-gas analyses suggested less methane and hydrogen and more carbon dioxide (\chem{CO_2}) and nitrogen. Repeating Miller-style experiments with these gases reduced yields, revising one premise of the original experiment. Scientists did not stop there. They changed mixtures and energy sources, including ultraviolet light, shock waves, and hydrothermal mineral surfaces, and found that amino acids and other parts could form under many plausible early conditions. Meteorites supplied another route: more than seventy amino acids were detected in the Murchison meteorite, which fell in Australia in 1969. Many were not used by terrestrial life and could not be explained by laboratory contamination. Organic parts could therefore form in space and might even have been ``delivered'' to early Earth by meteorites.

The point is not that ``Miller got the right answer,'' but the method: propose a hypothesis, challenge its premises, improve experiments, and open new channels of evidence. Seventy years later, ``Where did the parts come from?'' is largely answered, but ``How did they assemble into a replicating system?'' is not. That is science's normal breathing rhythm.
\end{zhengju}

\textbf{Translator safety note:} The spark-discharge apparatus is historical description, not a home experiment. Flammable gases, electrical ignition, and heated glass require a professionally assessed laboratory setup. See \href{https://www.energy.gov/cmei/fuels/safe-use-hydrogen}{DOE hydrogen precautions} and \href{https://institute.acs.org/acs-center/lab-safety/education-training/safer-experiments.html}{ACS experiment risk assessment}.

\begin{shiyan}
\textbf{A Pasteur-style controlled experiment: can soup ``produce'' life by itself?}

\textbf{Translator safety note:} Do not perform this unknown-microbe culture at home. The source's final directions to keep containers sealed but then scald and reuse them conflict; neither household disposal nor hot-water rinsing establishes safe decontamination. Use published observations, or an instructor-approved teaching laboratory with suitable containment and waste procedures. Never open, smell, taste, or attempt to clean an unknown culture yourself. See \href{https://asm.org/guideline/asm-guidelines-for-biosafety-in-teaching-laborator}{ASM teaching-laboratory biosafety guidance}. The original protocol follows for faithful translation, not as an endorsed procedure.

Materials: two equal portions of cooled rice soup, or water in which a little vegetable has been boiled; two clean clear cups with lids; plastic wrap; and a marker.

Procedure: (1) Fill each cup seven-tenths full with soup. (2) Leave cup A without its lid, loosely covering it with plastic wrap pierced with small holes so air can enter and leave. Seal cup B with its lid. (3) Place them side by side indoors in a bright place away from direct strong sunlight. Observe through the walls at the same time daily for five to seven days, recording cloudiness, surface films, and color changes. Look only: \textbf{do not open the lids or bring your face close to smell}.

What to observe: which cup clouds first? Do they change in the same way? If the open cup clouds first, does that support ``cloudiness comes from something entering from outside'' or ``soup spontaneously generates life''? Recall Pasteur's swan-neck flask from Chapter 45: its neck admitted air but blocked microbes. Which variable did he change?

Safety: cloudy soup contains multiplying unknown microbes. Throughout the experiment, \textbf{do not open or taste}. At the end, discard the whole sealed cup in the trash, scald the cups with boiling water, and wash your hands. Do not use dairy products, whose decay smells worse, or add anything intended to ``promote growth.''
\end{shiyan}

\section{From Folding to Cells to Consciousness}

The second unanswered question hides in scale. This book's causal chain keeps ``climbing floors'': atoms form molecules, molecules form cells, cells form tissues, and tissues form brains. At each floor, however, we discover that we cannot really ``calculate the upstairs behavior from downstairs rules.'' Several stairs are broken.

Start with proteins. Chapter 48 explained that a protein is an amino-acid chain that must fold into a precise three-dimensional shape to work. Here lies a famous numerical puzzle, \emph{Levinthal's paradox}.

\begin{qiaoliang}
Take a short protein of 100 amino acids. Suppose each amino acid has just 3 possible arrangements in the chain, far fewer than in reality. The whole chain then has approximately $3^{100}$ possible shapes. Estimate this: $\lg 3^{100} = 100 \times 0.477 \approx 47.7$, or approximately $5\times 10^{47}$ shapes. If the protein blindly tried shapes to find the correct one, spending only $10^{-13}$ seconds on each, roughly a molecular-motion timescale, testing them all would take approximately $5\times 10^{34}$ seconds. The universe is only approximately $4\times 10^{17}$ seconds old. In other words, blind trial would take more than a hundred million times a hundred million times a hundred million times a hundred million lifetimes of the universe. Yet real proteins in your cells fold in a millisecond to a few seconds: \textbf{a difference of more than thirty orders of magnitude}. Folding cannot be random search; it must follow paths guided by an energy landscape. What do those paths look like, and can they be calculated directly from sequence? That is the protein-folding problem. In recent years, AI programs such as AlphaFold have \emph{predicted} folded static structures quite accurately from sequences, greatly helping structural biology and drug design. But the \emph{process} of a chain folding inside a cell, assisted by molecular chaperones while still being synthesized, remains far from solved. Predicting the destination does not map the route.
\end{qiaoliang}

Above folding lie steeper stairs. Thousands of proteins, lipids, and small molecules react simultaneously within a cell, in Chapter 55's metabolic network. How do these reactions give rise to behavior such as a cell ``deciding'' whether to divide or wait? Chapter 57 only began the discussion. Higher still, your brain contains approximately 86 billion neurons, each obeying the chemistry of ion channels and neurotransmitters from Chapter 58. Yet together they produce the fact that ``you are reading and understanding this sentence.'' How does consciousness---subjective experience itself---arise from organized matter? We lack not only an answer but agreement on which quantities should measure it.

How difficult is this? Even the criteria are missing. Anesthetists can reliably make you lose consciousness and wake safely through experience and pharmacology, yet there is no unified explanation of exactly what anesthetics switch off. Doctors assess residual awareness in patients in coma using several indirect brain measurements whose underlying theories still compete. At the molecular level, we recognize the components; at the level of experience, we are still constructing the ruler. What lies between these levels is the broken staircase.

This does not mean the intervening floors are forever unknowable. It means that \textbf{predicting whole-system behavior from lower-level rules is among science's hardest jobs today}. The book's second recurring question---``How are the components connected and arranged?''---takes its most difficult form here.

\section{The Genome's Dark Map}

The third question returns to DNA. Your genome has approximately 3 billion letters, yet only approximately 1.5\% directly encodes proteins: the ``main text'' of roughly twenty thousand genes. What is the remaining 98\% or more?

Chapters 70 and 71 have already explained that it is far from merely junk. There are promoters and enhancers that switch genes, regions producing regulatory RNAs, structural chromosome components, and many repeats of uncertain origin, some of them ``fossils'' of ancient viruses. But how much has a function, what those functions are, and how important they are remain a half-drawn map.

The map is being revised in public, which makes it especially worth watching. In 2012, the large international ENCODE project announced detectable ``biochemical activity,'' such as transcription into RNA, across more than 80\% of the human genome. Many media outlets immediately reported, ``Junk DNA no longer exists!'' Evolutionary biologists objected: transcribing a sequence may be occasional idle running of cellular machinery. \textbf{Detectable activity does not equal function}. To judge genuine function, ask Chapter 76's question: would changing this sequence affect the organism? Does natural selection care about it? Under this stricter criterion, functional sequence might comprise only one or two tenths of the genome. The debate is not fully settled today, and it illustrates this chapter's theme: a bold announcement is tested, scaled back, and revised by colleagues applying stricter criteria. Science earns its ``knowledge'' inch by inch.

\section{Harder Than Weather: Genotype to Phenotype}

You may have guessed the fourth question, for which all of Chapter 73 prepared the ground: if we know someone's entire genetic sequence, can we predict what kind of person they will become?

Only a little, and often only probabilistically. For a trait such as height, hundreds or thousands of sites each contribute a few millimeters, alongside nutrition and childhood experience; statistical models now explain a substantial part of population height variation. For complex diseases such as heart disease, diabetes, and depression, genes provide only ``somewhat above- or below-average risk,'' mixed with environment, lifestyle, and sheer developmental randomness. Chapter 85 repeatedly emphasized that risk is not diagnosis. For personality, talent, and life choices, genetic predictive power is weak enough to be largely negligible.

Why so difficult? The causal chain crosses too many floors with feedback: genes regulate proteins, proteins shape cells, cells build brains, and brains respond to environments. Parts of that environment---family, school, and friends---in turn affect gene switching (Chapter 74). Identical twins provide especially strong evidence: their DNA sequences are nearly identical, and they often grow up under the same roof, yet one may have allergies and the other not, one be introverted and the other outgoing, or they may develop different diseases. Prediction fails even under the favorable conditions of nearly identical sequences and similar environments, showing how much chance enters development. A tiny error in one cell division or a neuron making one particular connection can leave lifelong marks. This resembles weather: too many entangled variables and sensitivity to initial conditions. People are even trickier than the atmosphere: atmospheric molecules do not read newspapers, but you do. Learning a prediction about yourself can change your behavior and invalidate it. Calculating phenotype from the joint effects of genotype, environment, and development is one of biology's most sought-after predictive achievements, and increasingly an ethical question everyone must face.

\section{Reverse the Question: Design the Future}

The first four questions concern understanding nature. The fifth reverses the book's direction: \textbf{prediction makes design possible}. Can we design sustainable materials, medicines, energy, and agricultural systems?

Clues to these tasks have appeared throughout the book. Energy: plants use chlorophyll and sunlight to fix carbon dioxide (Chapter 54), inefficiently but with the advantages of zero emissions and self-replication. Artificial photosynthesis aims to make fuels directly from sunlight, water, and \chem{CO_2} with catalysts. Finding inexpensive, durable catalysts is the challenge, making Chapter 30's knowledge valuable. Agriculture: nitrogen fertilizer feeding nearly half humanity comes from the century-old Haber process, which consumes a few percent of global energy annually. Bacteria in legume root nodules do the same job at ordinary temperature and pressure. Learning to imitate or improve biological nitrogen fixation is one of agricultural chemistry's great prizes. Materials: one route out of plastic pollution (Chapter 43) is polymers microbes can dismantle after use, requiring a joint understanding of polymer chemistry and metabolic networks. Medicine: folding predictions combined with gene sequences (Chapter 83) make ``a drug tailored to one mutation'' seem realistic for the first time.

These problems share a structure: goals lie upstairs---clean energy, enough food, treatable disease---while tools lie downstairs---catalysts, sequences, and networks. Connecting the levels is precisely the approach this book teaches. It also honestly reminds us that technology brings not only ability but the spillover and dual-use risks of Chapter 84. Anyone who knows how to design must also ask whether they should.

\begin{wuqu}
``If science admits it still does not understand everything, all explanations must be about equally reasonable. Homeopathy and water memory might yet prove right.'' This confuses two kinds of not knowing. One is \textbf{not yet having an answer}: several hypotheses for life's specific origin remain viable, without enough evidence for a verdict. The other is \textbf{already having a negative answer}: the homeopathic claim that dilution beyond the point of any remaining molecules can still cure disease has repeatedly performed no better than placebo in high-quality trials. The claim that water remembers what was dissolved in it also conflicts with the body of physical-chemical evidence (Chapter 23). Unknown frontiers are not a back door for rejected claims. The standard remains: can the claim be tested, and what do tests show? ``We do not yet know how X happened'' never means ``Every story about X is equally credible.''
\end{wuqu}

\textbf{Translator safety note:} Do not replace proven treatment or delay medical assessment with homeopathy. Some products labeled homeopathic contain active substances and can cause harm or interactions. See \href{https://www.nccih.nih.gov/health/homeopathy}{NCCIH guidance on homeopathy}.

\section{Ask That Bottle of Seawater Again}

Return to the opening bottle and direct the book's five recurring questions at it:

\begin{itemize}[itemsep=2pt]
\item What is it made of? Water, salts, and millions of bacteria and viruses, most unnamed. \textbf{Open}.
\item How are the parts connected and arranged? How these microbes form communities, eat one another, and send chemical signals was only introduced in Chapter 86; much remains unknown. \textbf{Open}.
\item Which energy and probability rules determine whether it can change? This is fairly clear: the same thermodynamics and kinetics as in a beaker. \textbf{Largely closed}.
\item How fast does change happen, and how is it controlled? Models exist for ocean \chem{CO_2} uptake, acidification rates, and carbon-cycle feedbacks, but their precision is limited. \textbf{Half-open}.
\item How is it preserved, copied, regulated, and selected in life? Every organism in this drop carries DNA and performs replication, regulation, and evolution. Yet how this system \emph{first} started is this chapter's first question. \textbf{Open at the origin}.
\end{itemize}

The five questions are not a questionnaire to discard after an exam, but tools to carry with you. Whenever something new appears---news of a breakthrough therapy, a new material, or a ``scientists discover'' notification---ask them in order, and you are already ahead of most readers.

\section{The Book Closes; the Questions Remain Open}

Eighty-six chapters ago, the book began with a grain of kitchen salt and followed atoms, molecules, reactions, cells, genes, and evolution to ecosystems and gene editing. This chapter deliberately stops beneath five unlit lamps. That is not discouragement, but the book's most important takeaway: science's credibility never comes from ``We know everything,'' but from recognizing what it does not know and knowing how to turn unknowns into knowledge. Propose testable hypotheses, let evidence judge, and revise when wrong. The periodic table predicted gallium; Miller's flask had its premises revised; ENCODE's 80\% is being reassessed. Every evidence story in the book shares this breathing rhythm.

There are signposts for the road ahead. Forget symbols or tools? Consult Appendices A, B, and C. Stuck on a term? Appendix D is a dictionary. Next time you see ``Research shows,'' run through Appendix E's checklist. Stuck on an exercise? Appendix F offers hints and further reading. More important than any appendix are the five questions you already carry. They belong to no single book, but to anyone willing to keep asking.

Years from now, you may encounter one of this chapter's questions firsthand in a laboratory, a field, or a clinic---or you may not. Wherever you are, one thing remains true, and it is the final sentence this book wants to leave with you:

You, leaves, medicine tablets, and distant stars are made of the same elements. Life's wonder is not that it escapes chemical laws, but that ordinary molecules can organize themselves, preserve information from the past, and create a future still capable of change.

\section*{Exercises}

\begin{sikaoti}
\item Choose a recent science news story and apply the five recurring questions. Identify which it answers, which it does not, and which it only pretends to answer. Organize your findings in a five-column table.
\item For each hypothesis, ``Life began at deep-sea vents'' and ``Life began in alternating wet and dry shallow pools,'' design an observation or experiment that could distinguish them. You may consider Earth, Mars, or Jupiter's moons. What result would support one and weaken the other?
\item Levinthal's paradox shows that folding is not blind trial. Draw an energy landscape with folding progress horizontally and energy vertically. Show a bumpy downward channel representing guided folding, alongside a contrasting random-search line. Note that the horizontal axis really represents an astronomical number of conformations; this diagram is only a thinking aid.
\item Only approximately 1.5\% of your genome encodes proteins. One friend says, ``The other 98.5\% is junk and can be deleted.'' Another says, ``All the other 98.5\% is treasure; not one part can be lost.'' Using ENCODE's distinction between activity and function, write three sentences responding to both.
\item ``Weather forecasts are inaccurate, so scientists cannot be trusted.'' Refute this using the chapter's three kinds of not knowing. Identify which kind forecast errors represent and how it fundamentally differs from claims already rejected by evidence.
\item The book's last exercise: choose any nearby object---a cup of water, a leaf, a pill---and use the five recurring questions to make a ``knowns and unknowns'' list. Under each question, write one thing science knows and one it does not. You will discover that you can already connect anything to the chain from atoms to evolution.
\end{sikaoti}
